% =====================================================================
% notacion.tex
% Responsabilidad única: presentar la Parte I en palabras sencillas y
% reunir en una tabla la notación que se usa en los ejercicios 1 a 5,
% para que ningún símbolo aparezca sin haberse definido antes.
% =====================================================================
\noindent En esta primera parte resolvemos a mano cinco ejercicios con
datos pequeños e inventados (cinco o seis observaciones). La idea no es
que los datos sean realistas sino poder seguir cada cuenta con lápiz y
papel, para entender qué hace por dentro un programa como scikit-learn
cuando, en la Parte II, le pedimos que ajuste un modelo con cientos de
datos reales. Los cinco ejercicios recorren los cuatro modelos del taller:
la regresión lineal simple (una variable explica a otra con una recta), la
regresión lineal múltiple (varias variables a la vez), la regresión
logística (predecir la probabilidad de que algo pase, por ejemplo que una
máquina falle) y los árboles de decisión (predecir haciendo preguntas
sucesivas de sí o no).

Todos los modelos de este taller son \emph{supervisados}: tenemos ejemplos
en los que conocemos tanto los datos de entrada (las \emph{variables
explicativas} o \emph{predictores}, que llamamos $x$) como el resultado que
queremos predecir (la \emph{variable respuesta}, que llamamos $y$), y
buscamos una regla que, a partir de $x$, se equivoque lo menos posible al
predecir $y$. La Tabla~\ref{tab:notacion} reúne los símbolos que usamos;
cada uno se vuelve a explicar la primera vez que aparece en un ejercicio.

\begin{table}[H]
\centering
\small
\begin{tabular}{>{\raggedright\arraybackslash}p{3.2cm}>{\raggedright\arraybackslash}p{11.3cm}}
\toprule
Símbolo & Significado \\
\midrule
$n$ & Número de observaciones (filas de datos). \\
$i=1,\dots,n$ & Índice que recorre las observaciones; $x_i$ y $y_i$ son
los valores de la observación número $i$. \\
$\sum_{i=1}^{n} a_i$ & Suma $a_1+a_2+\dots+a_n$. Cuando no hay ambigüedad
escribimos solo $\sum_i a_i$. \\
$\bar x,\ \bar y$ & Promedios: $\bar x=\frac1n\sum_i x_i$ y
$\bar y=\frac1n\sum_i y_i$. \\
$\beta_0,\ \beta_1,\dots$ & Parámetros del modelo, los números que el
ajuste debe encontrar ($\beta_0$ es el intercepto y $\beta_1,\beta_2,\dots$
las pendientes o coeficientes). \\
$\hat\beta_0,\ \hat\beta_1$ & Los valores de los parámetros que encontramos
con los datos (el gorro $\hat{\ }$ indica \emph{estimado a partir de los
datos}). \\
$\hat y_i$ & Predicción del modelo para la observación $i$. \\
$e_i = y_i-\hat y_i$ & Residuo: lo que el modelo no logró explicar en la
observación $i$ (positivo si el modelo se quedó corto, negativo si se
pasó). \\
$\mathrm{SSE}=\sum_i e_i^2$ & Suma de errores al cuadrado (\emph{sum of
squared errors}). \\
$\mathrm{MSE}=\mathrm{SSE}/n$ & Error cuadrático medio: el promedio de los
errores al cuadrado. \\
$\mathrm{SST}=\sum_i (y_i-\bar y)^2$ & Variabilidad total de $y$ alrededor
de su promedio. \\
$R^2=1-\mathrm{SSE}/\mathrm{SST}$ & Coeficiente de determinación: fracción
de la variabilidad de $y$ que el modelo explica (entre $0$ y $1$ en
entrenamiento). \\
$\partial S/\partial \beta$ & Derivada parcial: cuánto cambia $S$ cuando
movemos solo $\beta$ y dejamos fijo lo demás. En un mínimo todas valen
cero. \\
$X,\ \beta,\ y$ (negrita implícita) & En el Ejercicio~2, matriz de datos
(una fila por observación, una columna por variable), vector de parámetros
y vector de respuestas. \\
$X^\top$ & Transpuesta de $X$: intercambia filas por columnas. \\
$p$ & Probabilidad estimada de que ocurra el evento ($Y=1$). \\
$\sigma(z)=\frac{1}{1+e^{-z}}$ & Función sigmoide: convierte cualquier
número real $z$ en un valor entre $0$ y $1$. \\
$\log$ & Logaritmo natural (base $e\approx2.718$). \\
$t$ & Umbral: se predice la clase $1$ cuando $p\ge t$. \\
$FP,\ FN$ & Falsos positivos (se predijo $1$ y era $0$) y falsos negativos
(se predijo $0$ y era $1$). \\
$\nabla J$ & Gradiente: el vector con todas las derivadas parciales de
$J$; apunta hacia donde $J$ crece más rápido. \\
$\eta$ & Tasa de aprendizaje: el tamaño del paso en descenso de gradiente. \\
$\lambda,\ C$ & Intensidad de la regularización ($\lambda$ grande o $C$
pequeño castigan más los coeficientes grandes). \\
$G$ & Impureza de Gini de un nodo de un árbol. \\
\bottomrule
\end{tabular}
\caption{Notación de la Parte I.}
\label{tab:notacion}
\end{table}
